\chapter{Mite}

\section*{Definition and Overview}
Mites are microscopic arachnids belonging to the subclass Acari, distributed globally across diverse ecological niches. While many species are free-living, several are of major medical significance to humans, acting as direct ectoparasites or potent environmental allergens. The most prominent medically important mites include Sarcoptes scabiei (the scabies mite), which burrows into the human epidermis causing intense pruritus, and Demodex species (follicle mites), which inhabit pilosebaceous units and are implicated in dermatological disorders like rosacea and blepharitis. Additionally, Dermatophagoides species (house dust mites) do not colonise human skin directly but produce potent digestive enzymes in their faecal pellets that serve as major airborne allergens, triggering allergic rhinitis, asthma, and atopic dermatitis. Clinically, managing mite-related conditions involves identifying the specific species, treating active infestations with topical or systemic scabicides, managing hypersensitivity reactions, and implementing environmental control measures.

\section*{Epidemiology}
Mite-associated diseases present a global public health challenge, affecting hundreds of millions of people annually across all socio-economic groups. Scabies, in particular, is classified by the World Health Organization (WHO) as a neglected tropical disease, with a global prevalence of over 200 million cases at any given time. In Australia, scabies is hyperendemic in some remote Aboriginal and Torres Strait Islander communities, particularly affecting young children, where prevalence rates can reach up to 30--50\%. House dust mite allergy is one of the most common hypersensitivity disorders worldwide; in Australia, it affects up to 20--30\% of the population, representing a leading trigger for chronic asthma and perennial allergic rhinitis. Demodex colonisation increases progressively with age, being present in nearly 100\% of individuals over the age of 70.

\section*{Aetiology and Pathophysiology}
The pathomechanisms of mite-related conditions are driven by active parasitism, tissue invasion, or hypersensitivity reactions to mite antigens.
\begin{itemize}
    \item \textbf{Epidermal burrowing (Sarcoptes scabiei):} Fertilised female mites secrete proteolytic enzymes to dissolve the stratum corneum, creating burrows where they deposit eggs and faeces, triggering a delayed type IV hypersensitivity reaction.
    \item \textbf{Pilosebaceous infestation (Demodex):} Mites feed on sebum and follicular epithelial cells, leading to physical ductal obstruction, microvascular damage, and acting as vectors for coexisting bacteria (e.g. Bacillus oleronius).
    \item \textbf{Dust mite allergen inhalation:} Inhaled proteins (specifically Der p 1 and Der p 2) from house dust mite faecal particles degrade tight junctions in the respiratory epithelium, triggering IgE-mediated type I hypersensitivity.
    \item \textbf{Secondary bacterial infection:} Scratching scabietic lesions breaches the skin barrier, facilitating bacterial entry (most commonly Streptococcus pyogenes and Staphylococcus aureus), which can lead to impetigo and systemic glomerulonephritis.
    \item \textbf{Crusted scabies (hyperinfestation):} Occurs in immunocompromised patients (e.g. HIV-positive or elderly individuals), where a failure of cellular immunity allows millions of mites to proliferate, causing severe hyperkeratosis.
\end{itemize}

\section*{Clinical Features}
The clinical presentation depends on whether the symptoms are driven by parasitic cutaneous infestation or systemic allergic sensitisation.
\begin{itemize}
    \item \textbf{Intense nocturnal pruritus:} Excoriating itching that worsens dramatically at night, characteristic of scabies infestation.
    \item \textbf{Cutaneous burrows and papules:} Fine, thread-like, greyish-white lines on the skin (burrows) and small, erythematous papules, classically located in the interdigital webs, wrists, axillae, periumbilical area, and genitalia.
    \item \textbf{Facial erythema and papulopustules:} Redness, burning, and acne-like lesions on the cheeks, nose, and forehead, associated with Demodex overgrowth (rosacea).
    \item \textbf{Ocular irritation and blepharitis:} Gritty sensation in the eyes, redness, crusting of the eyelids, and loss of eyelashes, typical of Demodex blepharitis.
    \item \textbf{Allergic respiratory symptoms:} Paroxysmal sneezing, watery rhinorrhoea, nasal congestion, itchy eyes, wheezing, and dyspnoea, triggered by dust mite exposure.
\end{itemize}

\section*{Differential Diagnosis}
Mite infestations and allergies must be differentiated from other dermatological and respiratory conditions.
\begin{itemize}
    \item \textbf{Atopic eczema:} A chronic inflammatory skin condition presenting with itchy, dry plaques, which may coexist with or be triggered by dust mites but lacks the classic burrowing and distribution of scabies.
    \item \textbf{Insect bites (fleas, bed bugs):} Presenting with localised, grouped pruritic wheals, which typically do not involve the interdigital spaces or present with progressive nocturnal itching.
    \item \textbf{Acne vulgaris:} Characterised by comedones (blackheads or whiteheads), which are absent in Demodex-induced papulopustular rosacea.
    \item \textbf{Perennial allergic rhinitis (non-mite):} Chronic nasal symptoms caused by other environmental allergens (such as pet dander or mould spores), distinguished by allergy testing.
    \item \textbf{Seborrheic dermatitis:} Scaling and redness on the scalp and face, distinguished from Demodex infestation by its greasy appearance and distribution.
\end{itemize}

\section*{When to Seek Medical Care}
Individuals with persistent, unexplained itching (especially at night or affecting multiple household members), chronic facial redness, or progressive allergy symptoms should consult a general practitioner or dermatologist. Immediate emergency services should be contacted for severe systemic complications.

\textbf{Call 000 (or local emergency services) for:}
\begin{itemize}
    \item \textbf{Anaphylaxis:} Severe allergic reaction (e.g. difficulty breathing, swelling of the tongue, or drop in blood pressure) to topical scabicide treatments or during severe asthma flare-ups.
    \item \textbf{Severe secondary bacterial infection:} High fever, rapid heart rate, confusion, or spreading red, painful skin (erysipelas or cellulitis) complicating scratched scabies lesions.
\end{itemize}

\section*{Diagnosis and Investigations}
Diagnosis involves clinical inspection, microscopic visualisation, or immunological testing, depending on the suspected mite-associated pathology.
\begin{itemize}
    \item \textbf{Skin scraping and microscopy:} The diagnostic standard for scabies, involving scraping a burrow with a scalpel blade under mineral oil to visualise mites, eggs, or faecal pellets (scybala) under light microscopy.
    \item \textbf{Dermoscopy:} Visualising the "delta wing" sign (the dark, triangular anterior portion of the mite) directly within a skin burrow using a hand-held dermatoscope.
    \item \textbf{Standardised skin biopsy or cyanoacrylate strip:} Sampling facial sebum to count Demodex density (diagnostic of overgrowth if density exceeds 5 mites per square centimetre).
    \item \textbf{Skin prick testing (SPT):} Applying dust mite allergen extract to the forearm to detect IgE-mediated wheal-and-flare reactions, confirming house dust mite allergy.
    \item \textbf{Specific IgE blood tests:} Measuring circulating antibodies to dust mite proteins, useful if skin prick testing is contraindicated or unavailable.
\end{itemize}

\section*{Management}
Management is tailored to the specific mite condition, utilising antiparasitic agents, antihistamines, or environmental interventions.
\begin{itemize}
    \item \textbf{Topical scabicides:} Applying permethrin 5\% cream overnight from the neck down, washed off after 8--14 hours, and repeated 7 days later; all close household contacts must be treated concurrently.
    \item \textbf{Systemic antiparasitics:} Oral ivermectin (200 micrograms/kg, repeated in 7--14 days) indicated for refractory scabies, institutional outbreaks, or crusted scabies.
    \item \textbf{Topical anti-demodectic therapies:} Topical ivermectin 1\% cream or metronidazole gel for facial rosacea, and tea tree oil wipes or formulations for Demodex blepharitis.
    \item \textbf{Pharmacological allergy control:} Non-sedating oral antihistamines, intranasal corticosteroid sprays, and allergen immunotherapy (desensitisation) for dust mite allergies.
    \item \textbf{Environmental decontamination:} Washing all bed linen, towels, and clothing used in the previous 48 hours in hot water (above 60 degrees Celsius) and drying in a hot dryer to kill scabies mites.
\end{itemize}

\section*{Complications}
\begin{itemize}
    \item \textbf{Secondary bacterial pyoderma:} Impetigo or ecthyma caused by Streptococcus or Staphylococcus infection of scratched skin.
    \item \textbf{Post-streptococcal glomerulonephritis:} A renal complication resulting from streptococcal skin infections secondary to untreated scabies, particularly in paediatric populations.
    \item \textbf{Rheumatic fever and heart disease:} Chronic cardiovascular sequelae linked to recurrent streptococcal pyoderma in endemic regions.
    \item \textbf{Severe asthma exacerbations:} Poorly controlled asthma attacks triggered by persistent exposure to house dust mite allergens.
    \item \textbf{Ophthalmic damage:} Corneal neovascularisation or scarring resulting from chronic, untreated Demodex blepharitis.
\end{itemize}

\section*{Prognosis and Outcomes}
The prognosis for mite-related conditions is generally excellent with correct diagnosis and adherence to treatment protocols. Scabies infestations are typically cured within 2 weeks of proper scabicide application, though itching (post-scabetic itch) can persist for up to 4 weeks due to ongoing allergic reaction to dead mite antigen and does not represent treatment failure. Allergic rhinitis and asthma triggered by dust mites are highly manageable, with many patients achieving significant long-term relief through allergen immunotherapy and environmental control. Demodex-associated dermatological conditions respond well to targeted topical therapies but may require long-term maintenance treatment to prevent recurrences.

\section*{Prevention and Self-Care}
\begin{itemize}
    \item \textbf{Wash bedding weekly:} Washing sheets and pillowcases in hot water (at least 60 degrees Celsius) to kill house dust mites.
    \item \textbf{Use allergen-barrier covers:} Encasing mattresses, pillows, and quilts in dust-mite-proof protective covers.
    \item \textbf{Reduce indoor humidity:} Keeping indoor humidity levels below 50\% using dehumidifiers or air conditioning to limit dust mite survival.
    \item \textbf{Avoid direct skin contact:} Preventing close body contact with individuals diagnosed with active scabies until they complete treatment.
    \item \textbf{Vacuum carpets regularly:} Using a vacuum cleaner fitted with a HEPA filter to remove mite particles from carpets and soft furnishings.
\end{itemize}

\section*{Sources and Further Reading}
The following reputable organisations provide clinical guidelines and research on mites, scabies, and allergy management:
\begin{itemize}
    \item Australasian Society of Clinical Immunology and Allergy. \textit{Information and management guidelines for dust mite allergy.} https://www.allergy.org.au/
    \item Healthdirect Australia. \textit{Scabies symptoms, treatment, and environmental control.} https://www.healthdirect.gov.au/
    \item World Health Organization. \textit{Scabies and other ectoparasitic infestations resource portal.} https://www.who.int/
    \item DermNet New Zealand. \textit{Clinical guide to scabies, Demodex, and other mite bites.} https://dermnetnz.org/
    \item Mayo Clinic. \textit{Diagnosis and treatment of scabies and dust allergies.} https://www.mayoclinic.org/
\end{itemize>
